\documentclass{plantillaPPT}

\titulo{7}{Optimización y Multiplicadores de Lagrange}

\begin{document}
\primeraDiapo

\begin{frame}{Multiplicadores de Lagrange}
\framesubtitle{Definición}

\begin{teorema}
    Sean $f:A\subseteq\mathbb{R}^n\to\mathbb{R}$ y $g:A\subseteq\mathbb{R}^n\to\mathbb{R}$ funciones de clase $C^1$. Si $\vec{x_0}\in A$ es tal que $g(\vec{x_0}) = 0$ y $\nabla g(\vec{x_0}) \neq 0$, entonces para que $f$ alcance un extremo local en $\vec{x_0}$ es necesario que $\nabla f \parallel \nabla g$
\end{teorema}
\vspace{0.5cm}
El teorema establece que los extremos relativos de una función $f$ sometidos a una condición $g(\vec{x}) =0$ se obtienen resolviendo el sistema
\vspace{-0.1cm}
\[
\left\{
\begin{aligned}
g(\vec{x}) &= 0 \\
\nabla f(\vec{x}) &= \lambda\nabla g(\vec{x})
\end{aligned}
\right.
\]

\end{frame}

\begin{frame}{Problema 1}
\framesubtitle{Práctica 7}
1. Hallar los valores máximo y mínimo absolutos de la función
\[
f(x,y) = xy(1-x^2-y^2)
\]
en el rectángulo $R = [0,1]\times[0,1]$
\end{frame}

\begin{frame}{Problema 2}
\framesubtitle{Práctica 7}
2. Si $T(x,y)$ grados es la temperatura en cualquier punto $(x,y)$ de disco circular $D: x^2+y^2\leq 1$ y
\[
T(x,y) = 2x^3+y^2-y
\]
determine los puntos más calientes y los más fríos del disco $D$ y la temperatura en esos puntos.
\end{frame}

\begin{frame}{Problema 3}
\framesubtitle{Práctica 7}

3. Calcule las distancias menor y mayor desde el origen a un punto el elipsoide
\[
9x^2+4y^2+z^2 = 36
\]

\end{frame}


\end{document}